% !TEX root = ../projeto-pesquisa.tex
% =====================================================================
%  TEMA
% =====================================================================
\secao{Tema}
Análise de padrões e inconsistência para identificação de imagens geradas por IA,
através de ELA e EXIF.

% =====================================================================
%  JUSTIFICATIVA
% =====================================================================
\secao{Justificativa}
O desenvolvimento deste projeto permite que os estudantes apliquem na prática
conhecimentos relacionados à tecnologia, análise de dados, segurança da
informação e inteligência artificial. Durante a pesquisa e construção do
aplicativo ou \textit{site}, os alunos poderão desenvolver habilidades como
investigação tecnológica, pensamento crítico, programação e análise de metadados
digitais.

Além disso, o projeto contribui para a formação acadêmica ao aproximar o
conhecimento teórico da aplicação prática voltada para problemas reais da
sociedade, como a disseminação de \textit{fake news} e conteúdos manipulados pela
inteligência artificial.

Para a faculdade, o projeto fortalece as atividades de extensão universitária,
promovendo a integração entre ensino, pesquisa e sociedade. A iniciativa também
incentiva o desenvolvimento de soluções tecnológicas inovadoras que podem gerar
impacto social positivo, além de estimular a produção de conhecimento e a
conscientização digital da comunidade.

